\subsection{window-\texorpdfstring{$6$}{6} 规范群、体积与冻结 escort 的有限尺度闭式}

定理 \ref{thm:xi-foldbin-gauge-representation-distribution-law-m6-spectrum}、定理 \ref{thm:xi-foldbin-parity-charge-split-exact-minrank}、定理 \ref{thm:xi-foldbin-center-solvable-fiber-criterion}、定理 \ref{thm:xi-window6-center-exact-splitting-bdry-cyclic} 与定理 \ref{thm:xi-gauge-volume-boltzmann-stirling-exponential} 已分别给出 window-\(6\) 的规范群分解、奇偶荷分裂正合列、中心判别、boundary 中心直和分裂以及规范体积密度的 Stirling 指数律。把这些接口在 \(m=6\) 的刚性直方图 \(2:8,3:4,4:9\) 上收束后，可进一步得到下列完全闭式。

\begin{theorem}[window-\texorpdfstring{$6$}{6} 规范群的中心、导出群与 boundary 直和分裂]\label{thm:derived-window6-gauge-center-derived-splitting}
沿用定理 \ref{thm:xi-foldbin-gauge-representation-distribution-law-m6-spectrum} 与定理 \ref{thm:xi-window6-center-exact-splitting-bdry-cyclic} 的记号，记
\[
G_6^{\mathrm{bin}}:=\Aut(\Fold_6^{\mathrm{bin}}),
\qquad
Z_{\mathrm{bdry}}:=(\ZZ/2\ZZ)^3_{\mathrm{bdry}}\hookrightarrow Z\!\bigl(G_6^{\mathrm{bin}}\bigr).
\]
则有
\[
Z\!\bigl(G_6^{\mathrm{bin}}\bigr)\cong (\ZZ/2\ZZ)^8,
\]
\[
\bigl[G_6^{\mathrm{bin}},G_6^{\mathrm{bin}}\bigr]\cong A_3^{\,4}\times A_4^{\,9},
\]
并且存在分裂短正合列
\[
1\longrightarrow A_3^{\,4}\times A_4^{\,9}
\longrightarrow G_6^{\mathrm{bin}}
\longrightarrow (\ZZ/2\ZZ)^{21}
\longrightarrow 1.
\]
此外，\(Z_{\mathrm{bdry}}\) 在中心与 Abel 化中都是直和因子：
\[
Z\!\bigl(G_6^{\mathrm{bin}}\bigr)\cong Z_{\mathrm{bdry}}\oplus (\ZZ/2\ZZ)^5,
\qquad
\bigl(G_6^{\mathrm{bin}}\bigr)^{\mathrm{ab}}\cong Z_{\mathrm{bdry}}\oplus (\ZZ/2\ZZ)^{18}.
\]
\leanverified{paper\_derived\_window6\_gauge\_center\_derived\_splitting}
\end{theorem}
\begin{proof}
由定理 \ref{thm:xi-foldbin-gauge-representation-distribution-law-m6-spectrum}，
\[
G_6^{\mathrm{bin}}\cong S_2^{\,8}\times S_3^{\,4}\times S_4^{\,9}.
\]
对有限直积群，中心与导出子群逐因子分解，因此
\[
Z\!\bigl(G_6^{\mathrm{bin}}\bigr)
\cong
Z(S_2)^8\times Z(S_3)^4\times Z(S_4)^9,
\]
\[
\bigl[G_6^{\mathrm{bin}},G_6^{\mathrm{bin}}\bigr]
\cong
[S_2,S_2]^8\times [S_3,S_3]^4\times [S_4,S_4]^9.
\]
再用
\[
Z(S_2)=S_2\cong \ZZ/2\ZZ,\qquad Z(S_3)=Z(S_4)=1,
\]
\[
[S_2,S_2]=1,\qquad [S_3,S_3]=A_3,\qquad [S_4,S_4]=A_4,
\]
便得到
\[
Z\!\bigl(G_6^{\mathrm{bin}}\bigr)\cong (\ZZ/2\ZZ)^8,
\qquad
\bigl[G_6^{\mathrm{bin}},G_6^{\mathrm{bin}}\bigr]\cong A_3^{\,4}\times A_4^{\,9}.
\]

对各因子分别有分裂短正合列
\[
1\to 1\to S_2\to \ZZ/2\ZZ\to 1,
\qquad
1\to A_3\to S_3\to \ZZ/2\ZZ\to 1,
\qquad
1\to A_4\to S_4\to \ZZ/2\ZZ\to 1,
\]
其中截面均可由任一换位给出。对它们取直积，即得
\[
1\longrightarrow A_3^{\,4}\times A_4^{\,9}
\longrightarrow G_6^{\mathrm{bin}}
\longrightarrow (\ZZ/2\ZZ)^{8+4+9}
\longrightarrow 1,
\]
也即所示分裂正合列。特别地，
\[
\bigl(G_6^{\mathrm{bin}}\bigr)^{\mathrm{ab}}\cong (\ZZ/2\ZZ)^{21}.
\]

最后，中心分解
\[
Z\!\bigl(G_6^{\mathrm{bin}}\bigr)\cong Z_{\mathrm{bdry}}\oplus (\ZZ/2\ZZ)^5
\]
正是定理 \ref{thm:xi-window6-center-exact-splitting-bdry-cyclic} 的结论。另一方面，由定理 \ref{thm:xi-foldbin-parity-charge-split-exact-minrank}，
\[
\bigl(G_6^{\mathrm{bin}}\bigr)^{\mathrm{ab}}
\cong
(\ZZ/2\ZZ)^{\{w\in X_6:\ d_6(w)\ge 2\}}.
\]
而 window-\(6\) 的直方图只有 \(2,3,4\) 三层，故全部 \(21\) 条纤维都满足 \(d_6(w)\ge 2\)；再由定理 \ref{thm:xi-window6-center-exact-splitting-bdry-cyclic}，其中 boundary 集合 \(X_6^{\mathrm{bdry}}\) 恰含 \(3\) 个坐标。于是 Abel 化按坐标分裂为
\[
(\ZZ/2\ZZ)^{X_6^{\mathrm{bdry}}}\oplus (\ZZ/2\ZZ)^{X_6\setminus X_6^{\mathrm{bdry}}}
\cong
Z_{\mathrm{bdry}}\oplus (\ZZ/2\ZZ)^{18}.
\]
证毕。
\end{proof}

\begin{theorem}[window-\texorpdfstring{$6$}{6} 的规范体积与折叠缺陷精确差额]\label{thm:derived-window6-gauge-volume-defect-identity}
沿用定理 \ref{thm:xi-gauge-volume-boltzmann-stirling-exponential} 的记号。则
\[
|G_6^{\mathrm{bin}}|=2^{39}3^{13},
\]
\[
g_6=\frac{39\log 2+13\log 3}{64},
\qquad
\kappa_6=\frac{88\log 2+12\log 3}{64},
\]
并且
\[
\kappa_6-g_6=\frac{49\log 2-\log 3}{64}.
\]
等价地，
\[
\frac{e^{64\kappa_6}}{|G_6^{\mathrm{bin}}|}=\frac{2^{49}}{3}.
\]
\leanverified{paper\_derived\_window6\_gauge\_volume\_defect\_identity}
\end{theorem}
\begin{proof}
由定理 \ref{thm:xi-foldbin-gauge-representation-distribution-law-m6-spectrum}，
\[
G_6^{\mathrm{bin}}\cong S_2^{\,8}\times S_3^{\,4}\times S_4^{\,9},
\]
故
\[
|G_6^{\mathrm{bin}}|=(2!)^8(3!)^4(4!)^9.
\]
直接化简得
\[
|G_6^{\mathrm{bin}}|
=
2^8\cdot (2\cdot 3)^4\cdot (2^3\cdot 3)^9
=
2^{39}3^{13}.
\]
因此
\[
g_6=\frac1{64}\log |G_6^{\mathrm{bin}}|
=
\frac{39\log 2+13\log 3}{64}.
\]

另一方面，由 window-\(6\) 的纤维直方图 \(2:8,3:4,4:9\)，有
\[
\kappa_6
=
\frac1{64}\Bigl(8\cdot 2\log 2+4\cdot 3\log 3+9\cdot 4\log 4\Bigr).
\]
再用 \(\log 4=2\log 2\)，得到
\[
\kappa_6
=
\frac{16\log 2+12\log 3+72\log 2}{64}
=
\frac{88\log 2+12\log 3}{64}.
\]
两式相减即得
\[
\kappa_6-g_6
=
\frac{(88-39)\log 2+(12-13)\log 3}{64}
=
\frac{49\log 2-\log 3}{64}.
\]
最后，
\[
64\kappa_6-\log |G_6^{\mathrm{bin}}|
=
49\log 2-\log 3
=
\log\!\left(\frac{2^{49}}{3}\right),
\]
指数化即得
\[
\frac{e^{64\kappa_6}}{|G_6^{\mathrm{bin}}|}=\frac{2^{49}}{3}.
\]
证毕。
\end{proof}

\leanverified{paper\_derived\_window6\_representation\_zeta\_boundary\_extension}
\begin{theorem}[window-\texorpdfstring{$6$}{6} 规范群的表示 \texorpdfstring{$\zeta$}{zeta} 函数与 boundary parity 延拓计数]\label{thm:derived-window6-representation-zeta-boundary-extension}
定义
\[
\zeta_{G_6^{\mathrm{bin}}}^{\mathrm{irr}}(s)
:=
\sum_{\pi\in \widehat{G_6^{\mathrm{bin}}}}(\dim\pi)^{-s}.
\]
则有闭式
\[
\zeta_{G_6^{\mathrm{bin}}}^{\mathrm{irr}}(s)
=
2^8\,(2+2^{-s})^4\,(2+2^{-s}+2\cdot 3^{-s})^9.
\]
从而：
\begin{enumerate}
  \item 不可约复表示总数为
  \[
  \#\widehat{G_6^{\mathrm{bin}}}=2^8\cdot 3^4\cdot 5^9.
  \]
  \item 一维特征总数为
  \[
  \#\Hom(G_6^{\mathrm{bin}},\{\pm 1\})=2^{21}.
  \]
  \item restriction 映射
  \[
  \Hom(G_6^{\mathrm{bin}},\{\pm 1\})
  \longrightarrow
  \Hom(Z_{\mathrm{bdry}},\{\pm 1\})
  \]
  为满射，且每个 boundary parity 赋值恰有 \(2^{18}\) 个全局一维特征延拓。
\end{enumerate}
\end{theorem}
\begin{proof}
由
\[
G_6^{\mathrm{bin}}\cong S_2^{\,8}\times S_3^{\,4}\times S_4^{\,9},
\]
以及有限直积群的不可约表示由各因子不可约表示外张量积而成，表示 \(\zeta\)-函数在直积下相乘，故
\[
\zeta_{G_6^{\mathrm{bin}}}^{\mathrm{irr}}(s)
=
\zeta_{S_2}^{\mathrm{irr}}(s)^8
\zeta_{S_3}^{\mathrm{irr}}(s)^4
\zeta_{S_4}^{\mathrm{irr}}(s)^9.
\]
再用
\[
\zeta_{S_2}^{\mathrm{irr}}(s)=1+1=2,
\]
\[
\zeta_{S_3}^{\mathrm{irr}}(s)=1+1+2^{-s}=2+2^{-s},
\]
\[
\zeta_{S_4}^{\mathrm{irr}}(s)=1+1+2^{-s}+3^{-s}+3^{-s}=2+2^{-s}+2\cdot 3^{-s},
\]
便得到所示闭式。

令 \(s=0\)，即可得到不可约复表示总数
\[
\#\widehat{G_6^{\mathrm{bin}}}
=
2^8\cdot 3^4\cdot 5^9.
\]
一维特征数则等于 Abel 化大小，因此由定理 \ref{thm:derived-window6-gauge-center-derived-splitting}，
\[
\#\Hom(G_6^{\mathrm{bin}},\{\pm 1\})
=
\left|\bigl(G_6^{\mathrm{bin}}\bigr)^{\mathrm{ab}}\right|
=
2^{21}.
\]

同样由定理 \ref{thm:derived-window6-gauge-center-derived-splitting}，
\[
\bigl(G_6^{\mathrm{bin}}\bigr)^{\mathrm{ab}}
\cong
Z_{\mathrm{bdry}}\oplus (\ZZ/2\ZZ)^{18}.
\]
对 \(\Hom(-,\{\pm 1\})\) 取对偶后得到
\[
\Hom(G_6^{\mathrm{bin}},\{\pm 1\})
\cong
\Hom(Z_{\mathrm{bdry}},\{\pm 1\})\times (\ZZ/2\ZZ)^{18}.
\]
因此 restriction 映射是第一坐标投影，故满射；并且每个 boundary parity 赋值的纤维大小恰为 \(2^{18}\)。证毕。
\end{proof}

\begin{theorem}[window-\texorpdfstring{$6$}{6} 最大纤维层的 escort 冻结速率]\label{thm:derived-window6-escort-freezing-max-layer-rate}
\leanverified{paper\_derived\_window6\_escort\_freezing\_max\_layer\_rate}
取 \(m=6\)，并令
\[
d_6(w):=\bigl|\Fold_6^{\mathrm{bin}}{}^{-1}(w)\bigr|
\qquad (w\in X_6).
\]
由推论 \ref{cor:fold-bin-m6-fiber-hist}，其纤维直方图为 \(2:8,3:4,4:9\)。对任意 \(a>0\) 定义
\[
S_a(6):=\sum_{w\in X_6} d_6(w)^a
=
8\cdot 2^a+4\cdot 3^a+9\cdot 4^a,
\]
\[
\pi_a(w):=\frac{d_6(w)^a}{S_a(6)},
\qquad
M_j(a):=\pi_a\bigl(\{w\in X_6:\ d_6(w)=j\}\bigr)
\qquad (j=2,3,4).
\]
则有精确公式
\[
M_4(a)=\frac{9}{9+4(3/4)^a+8(1/2)^a},
\]
\[
M_3(a)=\frac{4(3/4)^a}{9+4(3/4)^a+8(1/2)^a},
\]
\[
M_2(a)=\frac{8(1/2)^a}{9+4(3/4)^a+8(1/2)^a}.
\]
从而：
\begin{enumerate}
  \item 冻结误差满足
  \[
  1-M_4(a)=\frac{4}{9}\left(\frac34\right)^a+O\!\left(\left(\frac{9}{16}\right)^a\right)
  \qquad (a\to\infty).
  \]
  \item 次大层与最小层的质量比满足精确恒等式
  \[
  \frac{M_3(a)}{M_2(a)}=\frac12\left(\frac32\right)^a.
  \]
  因而冻结误差 \(1-M_4(a)=M_3(a)+M_2(a)\) 最终由 \(d_6=3\) 层主导。
  \item 在最大层上，escort 条件分布对一切 \(a>0\) 都恰为均匀分布：
  \[
  \pi_a\bigl(\,\cdot\mid d_6=4\bigr)
  =
  \mathrm{Unif}\bigl(\{w\in X_6:\ d_6(w)=4\}\bigr).
  \]
\end{enumerate}
\end{theorem}
\begin{proof}
由直方图直接有
\[
S_a(6)=8\cdot 2^a+4\cdot 3^a+9\cdot 4^a.
\]
因此
\[
M_4(a)=\frac{9\cdot 4^a}{8\cdot 2^a+4\cdot 3^a+9\cdot 4^a}
=
\frac{9}{9+4(3/4)^a+8(1/2)^a},
\]
其余两式同理。

令
\[
r_a:=\left(\frac34\right)^a,
\qquad
s_a:=\left(\frac12\right)^a.
\]
则
\[
1-M_4(a)=\frac{4r_a+8s_a}{9+4r_a+8s_a}.
\]
由于 \(r_a,s_a\to 0\)，可写成
\[
1-M_4(a)
=
\frac{4r_a+8s_a}{9}\cdot \frac{1}{1+(4r_a+8s_a)/9}
=
\frac{4}{9}r_a+O(r_a^2+s_a).
\]
又 \(r_a^2=(9/16)^a\) 且 \(s_a=(1/2)^a=O((9/16)^a)\)，故
\[
1-M_4(a)=\frac49\left(\frac34\right)^a+O\!\left(\left(\frac{9}{16}\right)^a\right).
\]

第二条由精确公式直接化简：
\[
\frac{M_3(a)}{M_2(a)}
=
\frac{4(3/4)^a}{8(1/2)^a}
=
\frac12\left(\frac32\right)^a.
\]
因此 \(\frac{M_3(a)}{M_2(a)}\to\infty\)，从而冻结误差 \(M_3(a)+M_2(a)\) 渐近上由 \(M_3(a)\) 主导。

最后，在集合 \(\{w\in X_6:\ d_6(w)=4\}\) 上，每个点的 escort 权重都等于 \(4^a/S_a(6)\)。对该层条件化之后，得到的条件分布必为 \(9\) 个最大纤维上的均匀分布。证毕。
\end{proof}
\begin{theorem}[window-\texorpdfstring{$6$}{6} boundary parity 直因子的规范细化]\label{thm:derived-window6-boundary-parity-direct-factor-refinement}
\leanverified{paper\_derived\_window6\_boundary\_parity\_direct\_factor\_refinement}
沿用定理 \ref{thm:derived-window6-gauge-center-derived-splitting} 与定理 \ref{thm:xi-window6-center-exact-splitting-bdry-cyclic} 的记号。记
\[
X_6^{\mathrm{bdry}}
=
\{\texttt{100001},\texttt{100101},\texttt{101001}\},
\]
并对每个 \(w\in X_6^{\mathrm{bdry}}\) 令 \(\tau_w\) 为对应 \(S_2\)-因子中的非平凡元。记
\[
Z_{\mathrm{bdry}}:=\langle \tau_w:\ w\in X_6^{\mathrm{bdry}}\rangle\cong (\ZZ/2\ZZ)^3.
\]
再把其余五个 \(S_2\)-因子与全部 \(S_3^4\times S_4^9\) 合并记为
\[
G_6^{\mathrm{int}}:=S_2^{\,5}\times S_3^{\,4}\times S_4^{\,9}.
\]
则相对于既定的 boundary 纤维集，存在规范直积同构
\[
\boxed{\;
G_6^{\mathrm{bin}}
\cong
Z_{\mathrm{bdry}}\times G_6^{\mathrm{int}}
\cong
(\ZZ/2\ZZ)^3\times S_2^{\,5}\times S_3^{\,4}\times S_4^{\,9}.\;}
\]
特别地，window-\(6\) 的三条 boundary sheet parity 在群论层面形成一个真正的中心直因子，而其余 \(18\) 个稳定类型全部由 \(G_6^{\mathrm{int}}\) 承载。
\end{theorem}
\begin{proof}
由定理 \ref{thm:xi-foldbin-gauge-representation-distribution-law-m6-spectrum}，
\[
G_6^{\mathrm{bin}}\cong S_2^{\,8}\times S_3^{\,4}\times S_4^{\,9}.
\]
再由定理 \ref{thm:xi-window6-center-exact-splitting-bdry-cyclic}，boundary 集合 \(X_6^{\mathrm{bdry}}\) 恰对应其中三个 \(d_6=2\) 的坐标。故可将上式重排为
\[
G_6^{\mathrm{bin}}\cong
\left(\prod_{w\in X_6^{\mathrm{bdry}}}S_2\right)
\times
S_2^{\,5}\times S_3^{\,4}\times S_4^{\,9}.
\]
对每个 boundary 因子，\(\tau_w\) 正是其唯一非平凡元，因此
\[
\prod_{w\in X_6^{\mathrm{bdry}}}S_2
=
Z_{\mathrm{bdry}}
\cong
(\ZZ/2\ZZ)^3.
\]
于是得到所示直积同构。又因每个 boundary \(S_2\)-因子都落在 \(Z(G_6^{\mathrm{bin}})\) 中，故该直因子同时是中心直因子。证毕。
\end{proof}

\begin{theorem}[boundary parity 八个中心扇区上的复群代数完全等型化]\label{thm:derived-window6-boundary-sector-groupalgebra-isotypy}
\leanverified{paper\_derived\_window6\_boundary\_sector\_groupalgebra\_isotypy}
沿用定理 \ref{thm:derived-window6-boundary-parity-direct-factor-refinement} 的记号，并记
\[
\widehat Z_{\mathrm{bdry}}:=\Hom(Z_{\mathrm{bdry}},\{\pm1\}).
\]
对每个 \(\chi\in \widehat Z_{\mathrm{bdry}}\)，定义中心幂等元
\[
e_\chi:=\frac{1}{|Z_{\mathrm{bdry}}|}
\sum_{z\in Z_{\mathrm{bdry}}}\chi(z^{-1})\,z
\in \CC[G_6^{\mathrm{bin}}].
\]
则有精确分解
\[
\boxed{\;
\CC[G_6^{\mathrm{bin}}]
\cong
\bigoplus_{\chi\in\widehat Z_{\mathrm{bdry}}}
e_\chi\,\CC[G_6^{\mathrm{bin}}]
\cong
\bigoplus_{\chi\in\widehat Z_{\mathrm{bdry}}}
\CC[G_6^{\mathrm{int}}].\;}
\]
特别地，每个 boundary parity 扇区 \(e_\chi\CC[G_6^{\mathrm{bin}}]\) 在下列三种意义下都与其余七个扇区完全等型：
\begin{enumerate}
  \item 作为有限维半单 \(\CC\)-代数；
  \item 作为 Wedderburn 块维数多重集；
  \item 作为 Plancherel 权重分布。
\end{enumerate}
并且每个扇区的复维数恰为
\[
\boxed{\;
\dim_{\CC}\bigl(e_\chi\CC[G_6^{\mathrm{bin}}]\bigr)
=
|G_6^{\mathrm{int}}|
=
(2!)^5(3!)^4(4!)^9
=
\frac{|G_6^{\mathrm{bin}}|}{8}.\;}
\]
\end{theorem}
\begin{proof}
由定理 \ref{thm:derived-window6-boundary-parity-direct-factor-refinement}，
\[
G_6^{\mathrm{bin}}\cong Z_{\mathrm{bdry}}\times G_6^{\mathrm{int}},
\]
故
\[
\CC[G_6^{\mathrm{bin}}]
\cong
\CC[Z_{\mathrm{bdry}}]\otimes_{\CC}\CC[G_6^{\mathrm{int}}].
\]
又因 \(Z_{\mathrm{bdry}}\) 为有限 Abel 群，其复群代数按角色 Fourier 分解为
\[
\CC[Z_{\mathrm{bdry}}]
\cong
\bigoplus_{\chi\in\widehat Z_{\mathrm{bdry}}}\CC e_\chi.
\]
张量上 \(\CC[G_6^{\mathrm{int}}]\) 即得
\[
\CC[G_6^{\mathrm{bin}}]
\cong
\bigoplus_{\chi\in\widehat Z_{\mathrm{bdry}}}
(\CC e_\chi\otimes \CC[G_6^{\mathrm{int}}])
\cong
\bigoplus_{\chi\in\widehat Z_{\mathrm{bdry}}}
\CC[G_6^{\mathrm{int}}],
\]
并且每个直和因子正对应于 \(e_\chi\CC[G_6^{\mathrm{bin}}]\)。

由于八个扇区都与同一个半单代数 \(\CC[G_6^{\mathrm{int}}]\) 同构，它们的 Wedderburn 分块结构、不可约表示维数多重集以及由平方维数决定的 Plancherel 权重分布都完全一致。最后，维数公式由
\[
\dim_{\CC}\CC[G_6^{\mathrm{int}}]=|G_6^{\mathrm{int}}|
\]
与 \(|Z_{\mathrm{bdry}}|=8\) 立即给出。证毕。
\end{proof}

\begin{theorem}[boundary parity 扇区上的不可约表示计数与维数矩完全均匀]\label{thm:derived-window6-boundary-sector-irrep-moment-uniformity}
沿用定理 \ref{thm:derived-window6-boundary-sector-groupalgebra-isotypy} 的记号。对每个 \(\chi\in\widehat Z_{\mathrm{bdry}}\)，记
\[
\widehat{G_6^{\mathrm{bin}}}_\chi
:=
\left\{
\pi\in\widehat{G_6^{\mathrm{bin}}}:
\pi(z)=\chi(z)\,\mathrm{Id}
\ \text{对一切 } z\in Z_{\mathrm{bdry}}
\right\}.
\]
则对每个 \(\chi\) 都有严格恒等式
\[
\boxed{\;
\#\widehat{G_6^{\mathrm{bin}}}_\chi
=
\#\widehat{G_6^{\mathrm{int}}}
=
2^5\cdot 3^4\cdot 5^9.\;}
\]
并且对任意整数 \(r\ge 0\)，维数矩
\[
M_r(\chi):=
\sum_{\pi\in\widehat{G_6^{\mathrm{bin}}}_\chi}(\dim\pi)^r
\]
与 \(\chi\) 完全无关。特别地，
\[
\boxed{\;
\sum_{\pi\in\widehat{G_6^{\mathrm{bin}}}_\chi}(\dim\pi)^2
=
|G_6^{\mathrm{int}}|
=
(2!)^5(3!)^4(4!)^9.\;}
\]
从而 Plancherel 测度在八个 boundary parity 扇区上严格等分，每个扇区质量都等于 \(1/8\)。
\end{theorem}
\begin{proof}
由定理 \ref{thm:derived-window6-boundary-sector-groupalgebra-isotypy}，
\[
e_\chi\CC[G_6^{\mathrm{bin}}]\cong \CC[G_6^{\mathrm{int}}]
\]
对每个 \(\chi\) 都成立。有限群的不可约复表示与复群代数的 Wedderburn 简单块一一对应，因此
\[
\widehat{G_6^{\mathrm{bin}}}_\chi
\longleftrightarrow
\widehat{G_6^{\mathrm{int}}}.
\]
再由
\[
G_6^{\mathrm{int}}\cong S_2^{\,5}\times S_3^{\,4}\times S_4^{\,9},
\]
以及直积群不可约表示数按因子相乘，得到
\[
\#\widehat{G_6^{\mathrm{int}}}
=
\bigl(\#\widehat{S_2}\bigr)^5
\bigl(\#\widehat{S_3}\bigr)^4
\bigl(\#\widehat{S_4}\bigr)^9
=
2^5\cdot 3^4\cdot 5^9.
\]
这即给出第一式。

由于各扇区代数都同构于同一个 \(\CC[G_6^{\mathrm{int}}]\)，一切只由 Wedderburn 块尺寸决定的统计量，包括全部维数矩 \(M_r(\chi)\)，都与 \(\chi\) 无关。取 \(r=2\) 时，Burnside 恒等式给出
\[
\sum_{\rho\in\widehat{G_6^{\mathrm{int}}}}(\dim\rho)^2
=
|G_6^{\mathrm{int}}|,
\]
故
\[
\sum_{\pi\in\widehat{G_6^{\mathrm{bin}}}_\chi}(\dim\pi)^2
=
|G_6^{\mathrm{int}}|.
\]
将其除以
\[
|G_6^{\mathrm{bin}}|
=
8\,|G_6^{\mathrm{int}}|
\]
便知每个扇区承载的 Plancherel 总质量都恰为 \(1/8\)。证毕。
\leanverified{paper\_derived\_window6\_boundary\_sector\_irrep\_moment\_uniformity}
\end{proof}

\begin{theorem}[window-\texorpdfstring{$6$}{6} 群胚代数的 Elliott 不变量与 boundary face 的唯一识别]\label{thm:derived-window6-groupoid-elliott-boundary-face}
沿用命题 \ref{prop:fold-bin-groupoid-algebra-wedderburn} 与定理 \ref{thm:xi-window6-center-exact-splitting-bdry-cyclic} 的记号，定义
\[
\mathcal A_6:=\CC[\mathcal R_6^{\mathrm{bin}}].
\]
则由 window-\(6\) 的刚性直方图 \(2:8,3:4,4:9\) 有
\[
\boxed{\;
\mathcal A_6
\cong
M_2(\CC)^{\oplus 8}\oplus M_3(\CC)^{\oplus 4}\oplus M_4(\CC)^{\oplus 9}.\;}
\]
再把三条 boundary 纤维对应的三个 \(M_2(\CC)\)-块记为
\[
I_{\mathrm{bdry}}\cong M_2(\CC)^{\oplus 3},
\]
其余各块之直和记为
\[
I_{\mathrm{cyc}}
\cong
M_2(\CC)^{\oplus 5}\oplus M_3(\CC)^{\oplus 4}\oplus M_4(\CC)^{\oplus 9}.
\]
则
\[
\mathcal A_6=I_{\mathrm{bdry}}\oplus I_{\mathrm{cyc}},
\]
并且成立：
\begin{enumerate}
  \item Elliott 不变量完全闭式为
  \[
  K_0(\mathcal A_6)\cong \ZZ^{21},
  \qquad
  K_0(\mathcal A_6)^+\cong \NN^{21},
  \qquad
  K_1(\mathcal A_6)=0,
  \]
  其单位元类满足
  \[
  \boxed{\;
  [1_{\mathcal A_6}]
  =
  (\underbrace{2,\dots,2}_{8},
  \underbrace{3,\dots,3}_{4},
  \underbrace{4,\dots,4}_{9}).\;}
  \]
  \item tracial simplex 满足
  \[
  T(\mathcal A_6)\cong \Delta^{20},
  \qquad
  T(I_{\mathrm{bdry}})\cong \Delta^{2},
  \qquad
  T(I_{\mathrm{cyc}})\cong \Delta^{17},
  \]
  并且
  \[
  \boxed{\;
  T(\mathcal A_6)=T(I_{\mathrm{bdry}})\ast T(I_{\mathrm{cyc}}).\;}
  \]
  \item \(I_{\mathrm{bdry}}\) 是 \(\mathcal A_6\) 中唯一同时满足下列条件的理想：
  \begin{enumerate}
    \item 全部不可约表示维数都等于 \(2\)；
    \item 中心维数等于 \(3\)；
    \item 其单位元类在 \(K_0(\mathcal A_6)\) 中正好是由三条 boundary 纤维坐标支撑的三维坐标子向量。
  \end{enumerate}
\end{enumerate}
\end{theorem}
\begin{proof}
由命题 \ref{prop:fold-bin-groupoid-algebra-wedderburn}，
\[
\mathcal A_6
\cong
\bigoplus_{w\in X_6}M_{d_6^{\mathrm{bin}}(w)}(\CC).
\]
再由 window-\(6\) 的纤维直方图 \(2:8,3:4,4:9\)，立得
\[
\mathcal A_6
\cong
M_2(\CC)^{\oplus 8}\oplus M_3(\CC)^{\oplus 4}\oplus M_4(\CC)^{\oplus 9}.
\]
其中由定理 \ref{thm:xi-window6-center-exact-splitting-bdry-cyclic} 指定的三条 boundary \(d=2\) 纤维恰切出三个特殊的 \(M_2(\CC)\)-块，它们的直和正是 \(I_{\mathrm{bdry}}\)；其余各块的直和即为 \(I_{\mathrm{cyc}}\)。

对任意有限维 \(C^\ast\)-代数
\[
\bigoplus_{j=1}^r M_{n_j}(\CC),
\]
标准 \(K\)-理论给出
\[
K_0\cong \ZZ^r,
\qquad
K_0^+\cong \NN^r,
\qquad
K_1=0,
\]
且单位元类恰为 \((n_1,\dots,n_r)\)。把
\[
(n_1,\dots,n_{21})=(2^{\times 8},3^{\times 4},4^{\times 9})
\]
代入，即得第一条。

对迹空间，有限维 \(C^\ast\)-代数的 tracial state space 正是各块标准归一化迹的凸包，因此 \(r\) 个矩阵块对应 \(\Delta^{r-1}\)。故
\[
T(\mathcal A_6)\cong \Delta^{20},
\qquad
T(I_{\mathrm{bdry}})\cong \Delta^2,
\qquad
T(I_{\mathrm{cyc}})\cong \Delta^{17}.
\]
又因
\[
\mathcal A_6=I_{\mathrm{bdry}}\oplus I_{\mathrm{cyc}},
\]
所以任一迹都是两侧迹的凸组合，极端点正是两侧极端迹的并，故总体迹单纯形恰为两者的 join。

最后证明唯一性。设 \(J\subset \mathcal A_6\) 为任一理想并满足第三条中的三个条件。有限维半单代数的理想必为某些 Wedderburn 块的直和。由“全部不可约表示维数都等于 \(2\)”可知，\(J\) 只能由 \(M_2(\CC)\)-块组成；由“中心维数等于 \(3\)”可知，恰有三个这样的块；再由其单位元类在 \(K_0(\mathcal A_6)\) 中的坐标支撑正好是三条 boundary 纤维，可知这三个块只能是 boundary 所对应的那三块。因此 \(J=I_{\mathrm{bdry}}\)。证毕。
\end{proof}
\leanverified{paper\_derived\_window6\_groupoid\_elliott\_boundary\_face}

